\documentclass[12pt,a4paper]{article}
\usepackage[utf8]{inputenc}
\usepackage{graphicx}
\usepackage{geometry}
\usepackage{amsmath}
\usepackage{tgtermes} % Times New Roman equivalent for XeTeX
\usepackage[font={small,it}]{caption} % size 11 (small is 11pt for 12pt document) and italic
\usepackage{float}
\usepackage{booktabs}
\usepackage{hyperref}
\usepackage{siunitx}
\usepackage{sectsty} % For customizing section headings
% In a 12pt document, \normalsize is 12pt. \bfseries makes it bold.
\sectionfont{\fontsize{12}{14}\bfseries\selectfont} 

\geometry{
  a4paper,
  left=25mm,
  right=25mm,
  top=25mm,
  bottom=25mm,
}

\begin{document}

% --- Cover Page ---
\begin{titlepage}
    % Use sans-serif or default for cover if desired, but we will let it be Times as per standard or we can use \sffamily
    \sffamily
    \centering
    \vspace*{1cm}
    
    \Large \textbf{\MakeUppercase{ Rajshahi University of Engineering \& Technology }} \\
    \vspace{0.5cm}
    \large \textbf{\MakeUppercase{ Department of Electrical and Electronic Engineering }} \\
    
    \vspace{2.5cm}
    \Large \textbf{LABORATORY REPORT} \\
    \vspace{1cm}
    
    \begin{tabular}{ll}
        \textbf{Course No:} & EEE 3102 \\
        \textbf{Course Title:} & Digital Signal Processing Sessional \\
    \end{tabular}
    
    \vspace{2cm}
    
    \textbf{Experiment No:} 05 \\
    \vspace{0.5cm}
    \textbf{Name of the Experiment:} \\
    \Large \textbf{ Test Experiment about triac and diac together } \\
    
    \vspace{2.5cm}
    
    \begin{minipage}[t]{0.4\textwidth}
        \begin{flushleft} \large
            \emph{Submitted by:}\\
            John Doe \\
            Roll: 1901000 \\
            Section: A, Group: 1
        \end{flushleft}
    \end{minipage}
    \hfill
    \begin{minipage}[t]{0.5\textwidth}
        \begin{flushright} \large
            \emph{Submitted to:} \\
            
            Dr. Jane Smith \\
            Professor \\
            \vspace{0.3cm}
            
            Mr. Bob Brown \\
            Lecturer \\
            
            
        \end{flushright}
    \end{minipage}

    \vfill
    \textbf{Date of Performance:} October 09, 2026 \\
    \textbf{Date of Submission:} October 16, 2026
    
\end{titlepage}

% --- Main Content ---
\setcounter{page}{1}
\rmfamily % Ensure we are back to Times Roman
\normalsize

\begin{center}
    \textbf{Experiment No:} 05 \\
    \vspace{0.2cm}
    \textbf{Name of the Experiment:} Test Experiment about triac and diac together
\end{center}
\vspace{0.5cm}


\section{Objectives}
\begin{itemize}

    \item To observe the bidirectional conduction characteristics of a TRIAC and the symmetric breakover behavior of a DIAC.

    \item To investigate the phase control principles by analyzing the firing angle modulation of a TRIAC triggered by a DIAC in an AC circuit.

    \item To measure and plot the voltage and current waveforms across the load to verify power regulation functionality.

\end{itemize}


\section{Introduction}
The TRIAC (Triode for Alternating Current) and DIAC (Diode for Alternating Current) are essential semiconductor devices utilized in AC power control applications, such as light dimmers and motor speed regulators. Unlike a SCR, which conducts in only one direction, a TRIAC is a three-terminal, five-layer semiconductor device that can conduct current in both directions when triggered. It effectively functions as two SCRs connected in anti-parallel, allowing for full-wave AC control. The gating of the TRIAC is sensitive to the voltage applied between the gate and main terminal 2, and it turns off only when the current through it falls below the holding current, typically at the zero-crossing point of the AC waveform [1].

The DIAC is a two-terminal, four-layer semiconductor device that exhibits a high breakover voltage in both forward and reverse directions. It remains non-conductive until the voltage across it exceeds its breakover threshold, at which point it triggers into a low-impedance state. In a combined DIAC-TRIAC circuit, the DIAC is often used in series with a resistor-capacitor (RC) network to provide a reliable trigger pulse to the TRIAC gate. Because the DIAC triggers symmetrically in both half-cycles of the AC input, it ensures consistent phase control of the TRIAC regardless of the polarity of the supply, thereby eliminating the asymmetry that might occur with unipolar triggering methods [4]. This combination allows for precise adjustment of the power delivered to the load by varying the firing angle through an RC time constant network.



\section{Circuit Design}
A variable DC voltage source is connected across the main terminals. A gate current is provided to trigger the device. The voltage is swept from negative to positive values.


\begin{figure}[H]
    \centering
    \includegraphics[width=0.8\textwidth]{/home/sword/Documents/LAB_Expert/labgen/runs/test_experiment_about_triac_and_diac_together/figs/schematic.png}
    \caption{Experimental Circuit Diagram}
\end{figure}


\section{Required Apparatus}
\begin{itemize}

    \item DC Power Supply (Variable) (Quantity: 1)

    \item Device Under Test (Quantity: 1)

    \item Resistor ($1 k\Omega$) (Quantity: 1)

    \item Multimeter (Quantity: 2)

\end{itemize}

\section{Procedure}
\begin{enumerate}

    \item Connect the circuit as per the experimental circuit diagram.

    \item Apply a constant gate current $I_G$.

    \item Vary the supply voltage from -15V to 15V.

    \item Record the voltage and current.

    \item Plot the characteristics.

\end{enumerate}

\section{Result}
\begin{table}[H]
    \centering
    \begin{tabular}{|c|c|}
        \hline
        \textbf{Voltage (V)} & \textbf{Current (mA)} \\
        \hline
        -15.0 & -759.3 \\
        -11.7 & -753.0 \\
        -8.4 & -7.6 \\
        -5.0 & -3484.0 \\
        -1.7 & -1.5 \\
        1.6 & 712.9 \\
        5.0 & 4.3 \\
        8.3 & 748.2 \\
        11.7 & 10.9 \\
        15.0 & 762.8 \\
        \hline
    \end{tabular}
    \caption{Simulated Data (Sample Points)}
\end{table}



\begin{figure}[H]
    \centering
    \includegraphics[width=0.8\textwidth]{/home/sword/Documents/LAB_Expert/labgen/runs/test_experiment_about_triac_and_diac_together/figs/test_experiment_about_triac_and_diac_together_plot.png}
    \caption{ Simulated Test Experiment about triac and diac together Curve }
\end{figure}



\section{Discussion}
The experimental setup utilized a DIAC to trigger the gate of the TRIAC in a resistive load circuit, allowing for the observation of phase-controlled AC waveforms. The voltage across the DIAC was monitored until it reached its breakover voltage, at which point it conducted and supplied the necessary gate current to switch on the TRIAC. The simulation results showed that the TRIAC remained in a blocking state until the DIAC triggered the gate, after which the load voltage waveform appeared as a chopped version of the AC sine wave. The firing angle was observed to vary according to the phase of the AC supply, confirming that the DIAC provided symmetrical triggering for both the positive and negative half-cycles. The current through the load was found to be proportional to the voltage, with conduction ceasing at the natural zero-crossing of the current, which turned off the TRIAC for the remainder of the cycle. This behavior validated the theoretical model of AC phase control using bidirectional thyristors.

\section{Conclusion}
The experiment successfully demonstrated the cooperative operation of a TRIAC and a DIAC in an AC power control circuit. The DIAC exhibited the expected symmetric breakover characteristic, triggering the TRIAC gate at the calculated phase angle for both polarities of the AC supply. The TRIAC responded by conducting in both half-cycles, confirming its bidirectional switching capability, and blocked current until the appropriate gate signal was received. The measured waveforms indicated that the average power delivered to the load could be effectively modulated by controlling the firing angle via the DIAC triggering mechanism. Minor deviations in the sharpness of the turn-on edge were attributed to the internal capacitance of the devices and the parasitic inductance of the circuit wiring. Overall, the results verified that the DIAC-TRIAC combination is a robust and reliable method for achieving symmetrical phase control in AC applications such as light dimming and motor speed regulation.

\section{References}
\begin{enumerate}

    \item Generated by LabGen Knowledge Hub API

    \item Ngspice Simulation Data.

\end{enumerate}

\end{document}